\section{Method}
\label{sec:method}

\subsection{Overview}
Given a question $q$, supplied tables $\mathcal{T}$, and an associated passage
pool $\mathcal{D}$, the method constructs a bounded knowledge graph specifically
for $q$. The graph is an intermediate reasoning structure rather than a
persistent database. The pipeline is
\[
q \rightarrow \text{row and passage retrieval}
\rightarrow \text{local KG construction}
\rightarrow \text{grounded text paths}
\rightarrow \text{path ranking}
\rightarrow \text{answer}.
\]
No program generation or execution is required by the primary method. Python
and typed numerical programs remain separate experimental variants.

\subsection{Query-Driven Context Retrieval}
A deterministic anchor first inspects the question and table schema. Its current
rules recognize ordinal constraints, exact numerical table constraints, and a
small set of entity-bearing columns. An ordinal such as \emph{second} can
therefore identify a bridge entity from a rank--entity row even when that entity
does not occur in the question.

Table retrieval forms a union of anchor-matched rows and BM25-ranked rows, with a
global row budget $k$. Original row indices are retained after filtering. This
is important because the same label can occur in several rows and because
provenance must refer to the source table rather than the filtered position.
Passage retrieval first applies BM25 to $q$, then adds at most $b$ passages using
bridge-entity queries and deterministic schema expansions. Duplicate passages
are removed by source identifier. The current defaults are $k=5$ and $b=3$ in
the benchmark runner; experiments record any override.

The benchmark setting supplies an oracle table and its associated passage pool.
Retrieval is therefore query-local within supplied context, not open-domain
table or web search.

\subsection{Schema-Preserving Online KG Construction}
For each question, the system creates a directed multigraph
$G_q=(V_q,E_q)$. A semantic edge is
\[
e=(h,r,t,\pi),
\]
where $h$ and $t$ are entity or literal labels, $r$ is a normalized relation,
and $\pi$ stores source type, source identifier, and optional row, column, and
header-path metadata. A multigraph retains separate provenance records when
several sources express the same fact.

\paragraph{Deterministic table edges.}
Dataset adapters convert a selected row into a dictionary with a row label and
qualified column names. The first nonempty field supplies the row subject
$u_i$; every other nonempty cell produces
\[
(u_i,\operatorname{norm}(H_j),v_{ij},\pi_{ij}).
\]
Hierarchical headers are qualified before graph construction. For example,
\texttt{2002 / Dividend} becomes \texttt{2002\_dividend}, preventing a
dividend from being confused with \texttt{2002\_high}.

Direct subject--value edges are convenient for entity lookup but do not by
themselves preserve which values co-occur in one row. The path-text reasoning
view therefore adds a record node $\rho_i$:
\[
u_i \xrightarrow{\mathrm{has\_record}} \rho_i, \qquad
\rho_i \xrightarrow{\operatorname{norm}(H_j)} v_{ij}.
\]
For two records of the same company, this representation prevents a year from
one row being paired with a revenue from another. Record identifiers are local
implementation nodes and do not alter the legacy executable entity interface.

In parallel, an audit trace stores
Table $\rightarrow$ Row $\rightarrow$ Cell $\rightarrow$ Column topology.
Thus the implementation exposes an entity-centric semantic graph, a
record-aware path view, and a complete structural trace without conflating their
roles.

\paragraph{Text edges and entity resolution.}
An LLM extracts semantic triples from each retrieved passage. A small set of
deterministic extractors preserves recurrent numerical and name facts when the
model response is malformed. Failed optional table enrichment falls back to
deterministic cells; a failed passage is logged and skipped. Entity resolution
then applies lexical normalization, configured aliases, structural name rules,
and conservative string similarity. Resolution events and rejected self-loops
are included in the trace.

The graph metadata distinguishes entity-oriented nodes, literal attributes, and
record structure operationally. It is not presented as a learned ontology or
a general semantic type system.

\subsection{Grounded Text-Path Generation}
The reasoning edge set $E_q^*$ is the union of semantic edges and row-record
edges. Candidate generation never invents an edge: every emitted path step
contains an edge identifier from $E_q^*$ together with direction and provenance.
Traversal may follow an edge forward or backward, but its stored factual
direction remains unchanged.

Seeds are graph nodes that match a question phrase, a bridge entity, or an
explicit table constraint. When no node is matched, endpoints of relations
sharing terms with the question provide fallback seeds. From each seed, bounded
search enumerates simple connected paths of at most three edges. Branching and
the total enumerated candidates are capped. The implementation additionally
creates a row bundle containing
$u_i\rightarrow\rho_i$ and the most question-relevant outgoing cell edges.
A row bundle is a connected evidence subgraph rather than a forced linear chain;
this allows a year constraint and requested value to remain attached to the same
record.

Each candidate is verbalized directly from its verified edges, for example
\[
\texttt{A --has\_record--> row:7 --year--> 2022}
\]
together with a sibling edge
\[
\texttt{row:7 --revenue--> 10M}.
\]
This differs from asking a language model to propose plausible relation names
from a list of nodes and relations.

\subsection{Text-Path Scoring}
All candidates pass a hard groundedness gate because generation only accepts
registered edge identifiers. The implemented heuristic score is
\[
\begin{split}
S(p)={}&G+Q_e+Q_r+Q_c+C+M+R+B+A\\
      &-0.35\max(|E_p|-1,0),
\end{split}
\]
where $G=1$ is grounded-edge validity; $Q_e$ and $Q_r$ are question overlap
with path nodes and relations, each capped at 2; and $Q_c$ is question-term
coverage, also capped at 2. The constraint term $C=2$ when the path includes
an anchored table value or allowed output candidate.

The cross-modal completeness term $M=1.5$ when a path uses both table and text
evidence. This is not called agreement: different sources may supply different
hops. Same-fact corroboration $A=1$ is awarded only when normalized
$(h,r,t)$ is supported by both table and text. Row consistency $R=1$ requires
at least two record-aware edges from one source row, and entity binding $B=1$
additionally requires the \texttt{has\_record} edge. These weights are fixed
heuristics rather than learned probabilities.

The top $N$ distinct paths are serialized with their score components, nodes,
edges, source identifiers, row indices, columns, and header paths. Supporting
edges are packed within a character budget, with selected path edges receiving
priority. Context packing does not change the path scores.

\subsection{Answer Production and Audit Record}
The answer model receives only the ranked paths, supporting KG edges, table
constraints, and the question. It reasons over this evidence as text and returns
a short answer span; it is explicitly instructed not to generate or execute a
program. The output is canonicalized to a resolved KG alias and, when a table
constraint defines an allowed candidate set, to the corresponding exact table
value.

Each prediction records the retrieved row count, passage retrieval decisions,
entity anchor, KG summary, ranked paths, score decomposition, and selected path
identifier. Empty graphs or an empty candidate set produce an explicit failure
record rather than an unsupported answer.

The trace establishes that path edges exist in the constructed graph. It does
not prove that LLM-extracted text triples are factually correct, that the
retriever found all required evidence, or that the top-ranked grounded path is
semantically sufficient for the question.
